\section{More Evaluation Results on Other Datasets}
\label{sec:appendix-more-eval}

To further verify the generality of the proposed framework, additional experiments were conducted on two challenging benchmarks: a subset of the LlamaTouch \citep{zhang2024llamatouch} dataset containing social media tasks, and the AndroidWorld \citep{rawles2024androidworld} environment with modified XML extraction. The 64 tasks selected from LlamaTouch are detailed in Table~\ref{tab:task-list-llamatouch}, while all 116 tasks from AndroidWorld are used according to the task list provided in Appendix~F of the original paper \citep{rawles2024androidworld}.

\begin{longtable}{|p{1.8cm}|p{10.5cm}|}
\caption{64 tasks across 5 apps selected from LlamaTouch dataset.}
\label{tab:task-list-llamatouch}\\
\hline
\textbf{App} & \textbf{Task} \\
\hline
Instagram & Open the Instagram app and follow an account named `artem\_chek'. \\
          & Open the Instagram app and view the notification list. \\
        & Open the Instagram app and like the first post in the page. \\
        & Open the Instagram app and navigate to your personal profile page. \\
        & Open the Instagram app, navigate to your personal profile page and view your following list. \\
        & Open the Instagram app, search for `\#travelgoals' and follow the hashtag. \\
        & Open the Instagram app and navigate to `Settings'. \\
        & Open the Instagram app, go to `Settings' and activate `Private Account'. \\
         & Open the Instagram app and edit your profile to add `Travel Enthusiast' to your bio. \\
         % \hline
        & Open the Instagram app and view the message list. \\
        & Open the Instagram app and follow the user of first post . \\
        & Open the Instagram app, navigate to your personal profile page and view your `followers' list. \\
        \hline
        & Open the Instagram app, navigate to your personal profile page, go to `Settings' and view the 'Blocked' list. \\
        & Open the Instagram app, navigate to your personal profile page, go to `Settings' and view the `Close Friends' list. \\
        & Open the Instagram app, navigate to your personal profile page, go to `Settings' and view the `Archive' list. \\
\hline
Pinterest & Search for `interior design ideas' on Pinterest. \\
        & Open Pinterest and search a user named `Wendy's Lookbook' and follow the user. \\
        & Navigate to the settings menu in Piniterest. \\
        & Create a new board named `DIY Crafts' on Pinterest. \\
        & Open Pinterest and copy the link of a pin from the `Home' page. \\
        & Open Pinterest and save a pin from the `Home' page to your Profile. \\
        & Open Pinterest and view the comments of a pin from the `Home' page. \\
        & Open Pinterest and view your own pin list. \\
        & Open Pinterest and open your own profile page. \\
        & Open Pinterest, open your own profile page and copy the shared link of yourself. \\
        & Open Pinterest, open your own profile page and view your following list. \\
        & Open Pinterst and refresh the `Home' page to see new pins by clicking the `Home' button again. \\
        & Open Pinterest and open one of your own board. \\
        & View `Updates' in the `Notifications' page in Pinterest.\\
        & View `Messages' in the `Notifications' page in Pinterest. \\
        & Filter `Updates' notifications by `Comments' in the 'Notifications' page in Pinterest. \\
        & Open Pinterest and set your profile visibility to `Private profile' in the settings menu. \\
        & Open Pinterest and search `DIY' in your own board list and open it. \\
        & Open the Pinterest app, navigate to the `Saved' page and sort your boards by `A to Z'. \\
        & Open the Pinterest app, navigate to the `Saved' page and view the `Privacy and data' in the settings menu. \\
        & Open the Pinterest app, navigate to the `Saved' page and view the `Home feed tuner' in the settings menu. \\
        & Navigate to the `Saved' page in the Pinterst app and view the `Social permissions' in the settings menu. \\
        & Navigate to the `Saved' page in the Pinterst app, open the settings menu and enable the `Filter comments on my Pins' button in the `Social permissions'. \\
        & Open the Pinterest app, navigate to the `Saved' page and view the `Account management' in the settings menu. \\
\hline
Reddit & Subscribe to the subreddit r/science and set community alerts to frequent on Reddit. \\
        & Find and display the hottest post from r/worldnews on Reddit. \\
        & Search for all posts related to SpaceX launches last month in the r/space subreddit on Reddit. \\
        & Open Reddit, search for the `technology' subreddit, and access it. \\
        & Open Reddit, search for the `technology' subreddit and aceess to Communities. \\
        & Open Reddit, search for the `technology' subreddit and aceess to Media. \\
        & Open Reddit, navigate to your Home feed and select Latest. \\
        & Open Reddit, navigate to the Communities. \\
        & Open Reddit, navigate to the Communities, search for openai and join. \\
        & Open Reddit, navigate to create a community. \\
        & Open Reddit, navigate to Inbox. \\
        & Open Reddit, navigate to Notifications setting. \\
        & Open Reddit, navigate to Chat. \\
        & Open Reddit, navigate to Chat and explore channels. \\
        \hline
        & Open Reddit, check my profile. \\
        & Open Reddit, go to setting. \\
        & Open Reddit, go to history. \\
\hline
X(Twitter) & Open the X app and check my notifications \\
        & Find the latest post of Openai on X. \\
        & Follow Yann LeCun on X. \\
        & Check my inbox on X. \\
        & Check latest trendings in X app. \\
        & Upload my avatar on X app using my latest picture on device. \\
        & Bookmark latest post from OpenAI on X app. \\
        & Check posts from my bookmark on X app. \\
\hline
\end{longtable}

\textbf{Evaluation on Social Media Applications.}  
Social media apps typically feature highly dynamic and frequently updated UIs, posing additional challenges for layout-aware partitioning and selective UI reduction. Following this observation, experiments were performed on 64 tasks from popular social platforms, including Instagram, X (Twitter), Reddit, and Pinterest, sampled from the LlamaTouch dataset. The evaluation results are summarized in Table~\ref{tab:socialmedia-results}.

\begin{table}[htbp]
  \centering
  \caption{Performance comparison on a subset of the LlamaTouch dataset (social media apps). The proposed framework maintains comparable success while substantially reducing UI exposure.}
  \label{tab:socialmedia-results}
  \resizebox{0.8\columnwidth}{!}{
  \begin{tabular}{lcc}
    \toprule
    \textbf{Method} & \textbf{Success Rate} & \textbf{Reduction Rate} \\
    \midrule
    GPT-4o (Baseline) & 70.31\% (45/64) & 0.00\% \\
    Ours (GPT-4o + Gemma-2-9B-Instruct) & 60.94\% (39/64) & 48.94\% \\
    \bottomrule
  \end{tabular}
  }
\end{table}

The results indicate that the proposed CORE framework completed only six fewer tasks than the full-UI GPT-4o baseline, while achieving a 48.94\% reduction in UI exposure. These findings confirm the effectiveness of CORE on dynamic and visually complex social media UIs.

\textbf{Evaluation on AndroidWorld.}  
The AndroidWorld environment was adapted by modifying its XML extraction pipeline to match the methodology described in the main text, where the system relies on preprocessing XML trees extracted via \texttt{uiautomator2}. Evaluation results are summarized in Table~\ref{tab:androidworld-results}.

\begin{table}[!htbp]
  \centering
  \caption{Performance on the AndroidWorld environment using the modified XML extraction pipeline.}
  \label{tab:androidworld-results}
  \resizebox{0.85\columnwidth}{!}{
  \begin{tabular}{lcc}
    \toprule
    \textbf{Method} & \textbf{Success Rate} & \textbf{Reduction Rate} \\
    \midrule
    Qwen2.5-Max (Baseline) & 35.34\% (41/116) & 0.00\% \\
    Ours (Qwen2.5-Max + Gemma-2-9B-Instruct) & 27.59\% (32/116) & 36.96\% \\
    \bottomrule
  \end{tabular}
  }
\end{table}

The results show that the CORE framework achieves a 36.96\% reduction in UI exposure while maintaining a reasonable task success rate (27.59\%), compared to 35.34\% for the full-UI Qwen2.5-Max baseline. This further demonstrates the applicability of CORE to different task environments.
